\documentclass[preprint,12pt]{elsarticle}

\usepackage{amsmath,amssymb,amsfonts,amsthm}
\usepackage{booktabs}
\usepackage{array}
\usepackage{algorithm}
\usepackage{algpseudocode}
\usepackage{enumitem}
\usepackage{geometry}
\usepackage{hyperref}
\hypersetup{hidelinks}
\usepackage{microtype}
\usepackage{xcolor}
\usepackage{mathtools}

\geometry{margin=1in}
\journal{Computer Methods in Applied Mechanics and Engineering}
\biboptions{numbers,sort&compress}

\newtheorem{definition}{Definition}
\newtheorem{proposition}{Proposition}
\newtheorem{theorem}{Theorem}
\newtheorem{remark}{Remark}
\newtheorem{assumption}{Assumption}

\newcommand{\R}{\mathbb{R}}
\newcommand{\Sph}{\mathbb{S}}
\newcommand{\Nc}{\mathcal{N}}
\newcommand{\Cc}{\mathcal{C}}
\newcommand{\Pc}{\mathcal{P}}
\newcommand{\Tc}{\mathcal{T}}
\newcommand{\Oc}{\mathcal{O}}
\newcommand{\Gc}{\mathcal{G}}
\newcommand{\eps}{\varepsilon}
\newcommand{\dd}{\mathrm{d}}
\newcommand{\dist}{\operatorname{dist}}
\newcommand{\argmin}{\operatorname*{arg\,min}}
\newcommand{\diag}{\operatorname{diag}}
\newcommand{\clip}{\operatorname{clip}}
\newcommand{\proj}{\operatorname{proj}}
\newcommand{\calg}{\textsc{CALG}}
\newcommand{\method}{\textsc{CALG-Contact}}

\begin{document}

\begin{frontmatter}

\title{Contact-Aligned Local Graphs for Subdivision-Free High-Order Surface-to-Surface Contact}

\author[inst1]{Anonymous Author(s)}
\address[inst1]{Manuscript prepared for submission to Computer Methods in Applied Mechanics and Engineering}

\begin{abstract}
Surface-to-surface contact remains a limiting component in large-deformation finite element analysis, robotics, digital manufacturing, and industrial simulation. Traditional node-to-surface and facet-pair algorithms are robust and inexpensive, but their contact normals, gaps, and tangential directions remain tied to a piecewise-linear boundary representation. Increasing accuracy by global mesh refinement is often unattractive because it multiplies the number of collision primitives, candidate pairs, and active contact constraints. This draft proposes a contact-aligned local graph formulation for high-order surface-to-surface contact on arbitrary open or closed triangulated surfaces. The method does not require a global parameterization, a manually specified cut graph, or a master-slave partition. Instead, each broad-phase candidate pair automatically generates a local contact frame. When a normal-cone graphability certificate is satisfied, the regular surface-to-surface closest-point problem on two curved patches is reduced from a coupled four-dimensional problem to a scalar two-dimensional gap-field problem on the contact plane. When the certificate fails, the algorithm falls back to a generic curved patch-patch closest-point solve and, when needed, conservative subdivision or interval bounds. The paper draft develops the mathematical formulation, the certificate and conditioning estimates, the curved-jet primitive used to obtain subdivision-free high-order accuracy, a symmetric contact-patch potential, and a planned numerical validation program. Numerical results are reported from the current executable snapshot in the accompanying repository (`results/exampleX.csv`).
\end{abstract}

\begin{keyword}
Surface-to-surface contact \sep high-order contact \sep local graph representation \sep contact detection \sep curved triangle patches \sep finite element contact \sep computational contact mechanics \sep contact potentials
\end{keyword}

\end{frontmatter}

\section{Introduction}
\label{sec:introduction}

Contact is one of the most persistent sources of nonlinearity and algorithmic fragility in computational mechanics. Even when the bulk response is modeled by mature finite element or isogeometric discretizations, the contact interface couples geometry, inequality constraints, active-set changes, sliding, frictional directions, and local resolution in a manner that is difficult to regularize without losing either accuracy or robustness. The classical finite element literature has developed node-to-surface, segment-to-segment, mortar, penalty, augmented Lagrangian, and semi-smooth Newton formulations that are effective in a wide range of engineering settings \cite{laursen2002,wriggers2006,puso2004mortar,puso2004friction}. More recent variational barrier formulations such as Incremental Potential Contact (IPC) have demonstrated that contact can be embedded in a robust optimization framework with strong non-penetration behavior for large-deformation dynamics \cite{li2020ipc}. Nevertheless, the geometric front end used to detect, parameterize, and integrate contact remains a central bottleneck.

A major source of difficulty is that most contact algorithms ultimately interact with a discrete surface representation. If the boundary is piecewise linear, the gap function is only piecewise smooth and the surface normal is discontinuous across edges. During sliding, a contact point may cross an element boundary and experience an abrupt change of contact normal or tangential frame. These discontinuities can generate force oscillations and resolution-dependent behavior even when the continuum contact model is smooth. Recent work on smooth surface representations has highlighted that contact artifacts in piecewise $C^0$ geometry can significantly affect forward and inverse simulations, and that smooth representations can reduce spurious tangential forces \cite{du2024smooth}. High-order finite element contact has similarly shown that high-order geometry and bases can improve accuracy, yet practical implementations often continue to use dense linear surface proxies for collision and force transfer \cite{ferguson2023highorder}. Thus there is a gap between high-order mechanical discretizations and the low-order, globally refined, or proxy-based surface geometry used by the contact search.

The most direct route to higher contact accuracy is to refine the contact mesh. If each triangle is subdivided $k$ times, the number of surface primitives grows like $4^k$. Broad-phase traversal, near-phase candidate pairs, continuous collision detection, and contact constraints all inherit this growth. The alternative explored in this draft is to keep the original primitive count but replace each linear triangle with a local curved jet carrying position, normal, curvature, and conservative error information. The resulting primitive can support higher-order contact gaps and smoother normals without globally increasing the number of collision elements. However, using curved patches naively appears to make the narrow phase more expensive: a closest-point query between two parametric surface patches requires minimizing
\begin{equation}
    \Psi(p,q)=\frac{1}{2}\|X_A(p)-X_B(q)\|^2,\qquad p\in U_A,\; q\in U_B,
    \label{eq:four_dimensional_closest}
\end{equation}
which is a coupled four-dimensional problem for two two-parameter surfaces. In continuous collision detection it becomes a five-dimensional space-time problem, a fact also emphasized in recent work on parametric-surface CCD \cite{chen2024tdib}.

This paper draft proposes to avoid this cost in regular surface-to-surface contact by constructing a contact-aligned local graph on demand. For every candidate pair generated by a standard three-dimensional broad phase, a contact frame is built from the preliminary closest points and normals. Let $n_c$ be the contact-frame normal and let $\Pi_c=n_c^\perp$ be the contact plane. If the normal cones of the two patches satisfy
\begin{equation}
    N_A(x)\cdot n_c\ge \mu>0,
    \qquad
    -N_B(y)\cdot n_c\ge \mu>0,
    \label{eq:intro_graphability}
\end{equation}
then both patches can be represented locally as height graphs over the same plane,
\begin{equation}
    X_A(\xi)=\xi+h_A(\xi)n_c,
    \qquad
    X_B(\xi)=\xi+h_B(\xi)n_c,
    \qquad \xi\in \Pi_c.
\end{equation}
The regular contact gap becomes the scalar field
\begin{equation}
    g(\xi)=h_B(\xi)-h_A(\xi),
    \label{eq:intro_gap}
\end{equation}
posed on the projected overlap region. Under the graphability certificate, the four-dimensional closest-point search in Eq.~\eqref{eq:four_dimensional_closest} is replaced by the two-dimensional minimization or quadrature of $g$. If the certificate fails, the algorithm does not fail; it simply reverts to a generic curved patch-patch solve and, if needed, a conservative fallback. The central design principle is therefore not to impose global smoothness, topology, or parameterization assumptions on the input surfaces. Instead, local geometric assumptions become certificates that enable a fast path.

The formulation is intentionally different from using a global surface parameterization. Boundary-controlled conformal parameterization methods such as Boundary First Flattening (BFF) provide high-quality planar maps for disk-like domains at low amortized cost, using boundary data, Poincare-Steklov operators, and a fixed Laplace factorization \cite{sawhney2018bff}. Such maps are valuable for geometry processing and can serve as optional local conditioners for persistent contact patches. However, a global parameterization is not a suitable prerequisite for general contact: arbitrary surfaces may be closed, open, high genus, cut by manufacturing features, or dynamically modified. The present method therefore does not require global flattening, manual cutting, or seam management. It uses the contact itself to define the local coordinate system.

The proposed method, called \method{} in this draft, makes four contributions.
\begin{enumerate}[leftmargin=*,label=(\roman*)]
    \item \textbf{Contact-aligned local graph formulation.} For each candidate pair, the algorithm automatically constructs a local contact plane and, when certified, expresses both curved patches as height graphs over this plane. This removes the need for a global object parameterization.
    \item \textbf{Normal-cone graphability certificate.} A computable normal-cone condition gives a sufficient criterion for reducing regular surface-to-surface contact from a four-dimensional closest-point problem to a two-dimensional gap-field problem. The same certificate controls the conditioning of the local projection.
    \item \textbf{Subdivision-free high-order contact primitives.} Each triangle is lifted to a curved jet carrying $X$, $N$, curvature, and error bounds. High-order contact gaps and smoother normals are obtained without globally subdividing the collision mesh.
    \item \textbf{Symmetric contact patch potential and fallback hierarchy.} The two-dimensional gap field supports contact patch quadrature with geometrically exact projection-area weights. When graphability fails, a curved patch-patch solve and conservative subdivision/interval fallback preserve generality for boundaries, corners, and degenerate configurations.
\end{enumerate}

This manuscript is a draft. The mathematical method and the numerical-example design are complete enough to guide implementation, but numerical values, tables, and figures are intentionally marked as pending. An accompanying implementation plan specifies a modular software architecture, benchmark sequence, verification tests, and a reserved industrial-scale scenario.

\section{Related works}
\label{sec:related}

\subsection{Classical finite element contact and mortar formulations}

Computational contact mechanics has a long history in nonlinear finite element analysis. Standard references cover penalty, Lagrange multiplier, augmented Lagrangian, active-set, and semi-smooth formulations for small and large deformation contact \cite{laursen2002,wriggers2006}. Early node-to-surface methods are simple and efficient, but they introduce a master-slave bias and may over-constrain or under-integrate the interface when meshes are nonmatching or contact patches slide across element boundaries. Segment-to-segment and mortar approaches address many of these limitations by integrating contact quantities over interacting surface segments rather than enforcing contact only at nodes \cite{puso2004mortar,puso2004friction}. Mortar schemes are particularly influential because they improve consistency for nonmatching meshes and large sliding, but they require careful projection, segmentation, quadrature, and multiplier-space design.

Recent unbiased and midplane-based contact formulations continue to revisit the geometric symmetry of surface-to-surface interaction. Sahu and Petrinic proposed a single-pass segment-to-segment method that evaluates contact tractions with respect to a midplane and avoids prescribing a master and slave surface \cite{sahu2025midplane}. This trend is aligned with the present work: a physically meaningful surface-to-surface method should avoid bias when possible. The difference is that \method{} constructs a local contact-aligned graph and a symmetric gap-patch potential on demand, rather than relying on a fixed master surface or a global mid-surface discretization.

\subsection{Barrier potentials, IPC, and geometric contact potentials}

Variational barrier methods have reshaped contact simulation by embedding non-penetration in an optimization energy. IPC formulates implicitly time-stepped large-deformation dynamics with a barrier contact potential and achieves intersection- and inversion-free trajectories under its assumptions \cite{li2020ipc}. Extensions and accelerations have expanded IPC's scope, including codimensional contact, adhesion, GPU implementations, and high-order finite element formulations. However, the discrete nature of typical IPC contact pairs can introduce resolution-dependent behavior and spurious forces if the geometric representation is not sufficiently smooth.

Geometric Contact Potential (GCP) was introduced to derive a contact potential from continuum, smooth or piecewise smooth surfaces, with requirements such as barrier behavior, locality, differentiable dependence on shape, and absence of rest-state forces \cite{huang2025gcp}. GCP is directly relevant to the proposed method because it separates the continuum notion of contact from a particular discrete primitive list. \method{} can be interpreted as a geometric front end for such a potential: it provides the local gap field, active interaction set, projection-area weight, and quadrature samples used to approximate a continuum contact patch. In contrast to a purely discrete pair potential, the local graph formulation treats regular face-face contact as a two-dimensional surface integral.

\subsection{Smooth surfaces and high-order contact geometry}

Piecewise-linear contact surfaces have discontinuous normals, which can cause nonphysical jumps in contact forces during sliding. Du et al. identified this issue for IPC-like contact and showed that smooth surface representations can reduce tangential force irregularities \cite{du2024smooth}. High-order IPC similarly demonstrates that high-order bases and curved meshes can improve accuracy and reliability in elastodynamic contact simulations \cite{ferguson2023highorder}. A key implementation compromise in existing high-order IPC is that collision detection and contact forces may still be evaluated on a dense linear surface proxy and transferred to the high-order representation.

The present method follows the smooth-contact motivation but changes the geometric interface. It keeps the number of broad-phase surface primitives close to the original mesh while enriching each primitive into a curved jet. The broad phase remains compatible with conventional three-dimensional BVHs, but the narrow phase and patch quadrature use curved geometry. This design aims to obtain the accuracy and normal continuity benefits of high-order geometry without the global candidate-count growth caused by dense collision proxies.

\subsection{Curved collision detection and conservative fallback}

Curved primitives improve geometric fidelity but complicate robust collision detection. Continuous collision detection between two parametric surfaces can be formulated as a five-dimensional constrained optimization over two surface parameters and time. Chen et al. developed a time-dependent inclusion-based method to reduce the cost of parametric-surface CCD by avoiding temporal subdivision \cite{chen2024tdib}. Sum-of-squares, Bernstein range bounds, convex hull subdivision, and interval Newton methods are also natural tools for certified queries on polynomial curves and surfaces.

\method{} uses these techniques as a fallback rather than as the default path. In regular surface-to-surface contact, the normal-cone graphability certificate reduces the local query to a two-dimensional gap field. Only when graphability fails, the local Newton solve becomes uncertain, or a CCD guarantee is required does the method invoke higher-cost conservative subdivision or interval machinery. This output-sensitive hierarchy is central to the proposed efficiency claim.

\subsection{Signed-distance and offset-based contact}

Signed distance fields provide continuous geometry and efficient queries, especially for rigid or volumetric collision environments. Recent triangle-SDF CCD formulations cast the first time of impact as a spatio-temporal local optimization problem \cite{pelletier2025trisdf}. Offset Geometric Contact constructs volumetric offset shapes from faces to obtain penetration-free behavior with orthogonal contact forces and low overhead for codimensional objects \cite{chen2025ogc}. These approaches show the value of continuous, locally optimizable contact geometry.

The method in this draft remains explicitly surface based. It does not require volumetric distance fields or global offset volumes, although it can incorporate an offset radius in the gap threshold and bounding boxes. This is advantageous for open shells, surface meshes extracted from CAD, and high-order finite element boundaries where rebuilding a volumetric field may introduce additional approximation error.

\subsection{Surface parameterization and the role of conformal maps}

Surface parameterization is an important tool in geometry processing and simulation. BFF provides boundary-controlled conformal flattening of disk-topology meshes with a workflow based on compatible boundary data, boundary curve integration, and harmonic or holomorphic extension \cite{sawhney2018bff}. Its computational appeal is that multiple Poisson-type problems share a fixed Laplace matrix. However, BFF and related conformal flattening methods are not used as mandatory contact coordinates in this work. Global parameterizations require topology management, cuts, seams, and chart overlap handling; these are undesirable prerequisites for an arbitrary contact algorithm. We instead use local contact-aligned coordinates generated per candidate pair. BFF or other parameterization tools may still be used as optional conditioners for persistent contact patches, for visualization, or for adaptive quadrature on large stable contact regions.

\section{Method}
\label{sec:method}

\subsection{Problem statement and design requirements}

Let two bodies have boundary surfaces $S_A,S_B\subset\R^3$. The input representation is a triangulated surface mesh, possibly open or closed, with no requirement of global parameterization or disk topology. The surface may contain smooth regions, sharp features, boundaries, and nonmatching discretizations. The method must return either contact samples for a downstream finite element or dynamics solver, or a local contact-patch energy suitable for penalty, augmented Lagrangian, or barrier enforcement.

The method is designed around the following requirements.
\begin{enumerate}[leftmargin=*]
    \item \textbf{Topological generality.} The algorithm must not require a global chart, user-specified cut, or seam handling.
    \item \textbf{Local smoothness when available.} Smooth curved geometry should be exploited in regular contact regions, but sharp or ambiguous regions must automatically fall back to generic queries.
    \item \textbf{Output-sensitive cost.} High-order operations should scale with the number of near or active contact pairs, not with a globally refined collision mesh.
    \item \textbf{Bias reduction.} Surface-to-surface contact should be formulated symmetrically whenever the local geometry permits.
    \item \textbf{Failure containment.} A failed graph certificate or Newton iteration should not invalidate the global contact search; it should only select a more general local solver.
\end{enumerate}

\subsection{Curved jet primitives}
\label{subsec:curvedjet}

Each input triangle $T_i$ is lifted into a curved jet primitive
\begin{equation}
    \Pc_i=\left(X_i,N_i,\kappa_i^{\max},\eps_i^x,\eps_i^N,\Cc_i^N\right),
\end{equation}
where $X_i:U_i\rightarrow\R^3$ is a local curved patch over a reference triangle $U_i$, $N_i$ is a smooth or piecewise smooth normal field, $\kappa_i^{\max}$ is a conservative curvature bound, $\eps_i^x$ and $\eps_i^N$ are position and normal error bounds relative to the intended geometry, and $\Cc_i^N$ is a normal cone enclosing $N_i(U_i)$.

The draft implementation uses either a cubic Bezier triangle or a quadratic normal-height primitive. In the Bezier case,
\begin{equation}
    X_i(u,v,w)=\sum_{a+b+c=d}B_{abc}^{d}(u,v,w)c_{abc},
    \qquad d=3,
    \qquad u+v+w=1,
    \label{eq:bezier_patch}
\end{equation}
where $B_{abc}^{d}$ are Bernstein basis functions. Vertex control points coincide with the original triangle vertices, while edge and interior control points are chosen from vertex normals or a local least-squares fit. This construction is related to curved PN triangles \cite{vlachos2001pn}. The normal field may be computed analytically from the patch,
\begin{equation}
    N_i=\frac{X_{i,u}\times X_{i,v}}{\|X_{i,u}\times X_{i,v}\|},
\end{equation}
or represented by a separate fitted normal field with consistency checks.

For high-order accuracy, the intended theoretical assumption is the following.
\begin{assumption}[Curved-jet approximation]
\label{assump:jet_accuracy}
Let $S$ be locally $C^3$ and sampled with characteristic edge length $h$. The curved jet primitive approximates $S$ with
\begin{equation}
    \|X_i-X\|_{L^\infty(U_i)}\le C_x h^3,
    \qquad
    \angle(N_i,N)\le C_N h^2,
    \label{eq:jet_accuracy}
\end{equation}
away from intentionally sharp features. In implementation, the constants are replaced by computable conservative bounds $\eps_i^x$ and $\eps_i^N$.
\end{assumption}

Assumption~\ref{assump:jet_accuracy} is not required for robustness; it is used to state convergence expectations. If no reliable smooth surface estimate is available, the primitive can be downgraded to the original triangle with larger error bounds.

\subsection{Conservative three-dimensional broad phase}

The global candidate search is purely three-dimensional. Each primitive receives an inflated bounding box
\begin{equation}
    \operatorname{AABB}_i^{\rm inf}
    =\operatorname{AABB}(X_i(U_i))\oplus
    \left(\eps_i^x+\hat d+r_i+\eps_i^{\rm motion}\right),
    \label{eq:inflated_aabb}
\end{equation}
where $\hat d$ is the contact activation distance, $r_i$ is an optional physical shell or offset radius, and $\eps_i^{\rm motion}$ is a swept bound for a time step. BVH traversal produces candidate pairs. A linear triangle distance is used as a first rejection test:
\begin{equation}
    d_{\rm lin}(T_i,T_j)>\hat d+\eps_i^x+\eps_j^x+r_i+r_j
    \quad\Longrightarrow\quad \text{reject}.
\end{equation}
Pairs that remain are passed to the contact-aligned narrow phase.

\subsection{Automatic contact-aligned frame}
\label{subsec:contact_frame}

For a candidate pair $(\Pc_A,\Pc_B)$, compute preliminary closest points $(x_A^0,x_B^0)$ using the linear triangles or a low-cost curved query. Define
\begin{equation}
    x_c=\frac{1}{2}(x_A^0+x_B^0).
\end{equation}
If $\|x_B^0-x_A^0\|$ is sufficiently large, set
\begin{equation}
    n_c=\frac{x_B^0-x_A^0}{\|x_B^0-x_A^0\|}.
    \label{eq:contact_normal_from_gap}
\end{equation}
If the preliminary gap is near zero, use the opposing normals,
\begin{equation}
    n_c=\frac{\bar N_A-\bar N_B}{\|\bar N_A-\bar N_B\|},
    \label{eq:contact_normal_from_normals}
\end{equation}
where $\bar N_A$ and $\bar N_B$ are representative patch normals. A deterministic rule chooses an orthonormal basis $(t_1,t_2,n_c)$.

The contact-plane coordinates of a point $x$ are
\begin{equation}
    \xi(x)=
    \begin{bmatrix}
    t_1^T(x-x_c)\\[2pt]
    t_2^T(x-x_c)
    \end{bmatrix},
    \qquad
    h(x)=n_c^T(x-x_c).
    \label{eq:projection_height}
\end{equation}
The contact plane is not a global parameterization of either object. It is created only for this local pair and discarded or cached after the query.

\subsection{Normal-cone graphability certificate}
\label{subsec:certificate}

Let $\pi_c(x)=\xi(x)$ be projection onto the contact plane. A patch is graphable over $\Pi_c$ if $\pi_c$ is locally one-to-one and nonsingular on the patch. A sufficient condition is that the patch normal never becomes tangent to the contact direction.

\begin{definition}[Normal-cone graphability]
A pair $(\Pc_A,\Pc_B)$ is $\mu$-graphable in the contact frame $(t_1,t_2,n_c)$ if
\begin{equation}
    \inf_{p\in U_A} N_A(p)\cdot n_c\ge \mu>0,
    \qquad
    \inf_{q\in U_B} -N_B(q)\cdot n_c\ge \mu>0.
    \label{eq:mu_graphability}
\end{equation}
\end{definition}

The condition is evaluated conservatively by normal cones. If the normal cone of patch $A$ has axis $\bar N_A$ and angular radius $\theta_A$, then
\begin{equation}
    \angle(\bar N_A,n_c)+\theta_A < \frac{\pi}{2}-\delta
    \label{eq:cone_test_a}
\end{equation}
ensures $N_A\cdot n_c\ge \mu=\sin\delta$. A similar test is applied to $-N_B$.

\begin{theorem}[Certified local graph reduction]
\label{thm:graph_reduction}
Let $X:U\subset\R^2\rightarrow\R^3$ be a $C^1$ regular patch with unit normal $N$. If $N(p)\cdot n_c\ge \mu>0$ for all $p\in U$, then the orthogonal projection $\pi_c\circ X$ is locally invertible everywhere on $U$. Moreover, the inverse graph map has condition number bounded by
\begin{equation}
    \|D(\pi_c\circ X)^{-1}\|\le \frac{C_X}{\mu},
    \label{eq:inverse_condition_bound}
\end{equation}
where $C_X$ depends only on the local parameter scaling of $X$. Consequently $X$ can be locally written as a height graph $X(\xi)=\xi+h(\xi)n_c$ over the contact plane. If both contacting patches are $\mu$-graphable over the same plane and have overlapping projections, their regular surface-to-surface gap is represented by the scalar field $g(\xi)=h_B(\xi)-h_A(\xi)$ on the overlap.
\end{theorem}

\begin{proof}[Sketch]
The Jacobian of $\pi_c\circ X$ loses rank exactly when a nonzero tangent vector of the surface projects to zero, i.e., when that tangent vector is parallel to $n_c$. This occurs only if $n_c$ lies in the tangent plane, equivalently $N\cdot n_c=0$. The lower bound $N\cdot n_c\ge\mu$ therefore prevents singular projection. Standard inverse function arguments give local graphability and an inverse norm controlled by the angle between $n_c$ and the tangent plane. Applying the construction to both patches over a shared contact plane yields the common gap field.
\end{proof}

\begin{remark}
The graphability certificate is a fast-path condition, not an input requirement. If it fails, the algorithm uses the generic fallback in Sec.~\ref{subsec:patch_pair_fallback}.
\end{remark}

\subsection{Two-dimensional gap field and contact patch}
\label{subsec:graph_gap}

Assume $(\Pc_A,\Pc_B)$ is $\mu$-graphable. Let $D_A=\pi_c(X_A(U_A))$ and $D_B=\pi_c(X_B(U_B))$. The projected interaction domain is
\begin{equation}
    D_{AB}=D_A\cap D_B.
\end{equation}
For a point $\xi\in D_{AB}$, let $p_A(\xi)$ and $p_B(\xi)$ be the inverse projected coordinates on the two patches. The height functions are
\begin{equation}
    h_A(\xi)=n_c^T(X_A(p_A(\xi))-x_c),
    \qquad
    h_B(\xi)=n_c^T(X_B(p_B(\xi))-x_c),
\end{equation}
with signed gap
\begin{equation}
    g_{AB}(\xi)=h_B(\xi)-h_A(\xi)-r_A-r_B.
    \label{eq:graph_gap}
\end{equation}
Contact is active where $g_{AB}(\xi)<\hat d$. The simplest detector finds local minima of $g_{AB}$ in $D_{AB}$; the patch-energy version integrates a contact potential over the active part of $D_{AB}$.

For a graph surface with unit normal $N$ and projection direction $n_c$, the area element satisfies
\begin{equation}
    \dd A = \frac{1}{|N\cdot n_c|}\,\dd \xi.
    \label{eq:area_weight}
\end{equation}
Therefore a symmetric projection-area weight is
\begin{equation}
    w_{AB}(\xi)=\frac{1}{2}
    \left(
    \frac{1}{|N_A(p_A(\xi))\cdot n_c|}
    +
    \frac{1}{|N_B(p_B(\xi))\cdot n_c|}
    \right).
    \label{eq:symmetric_weight}
\end{equation}
The local contact patch potential is defined as
\begin{equation}
    E_{AB}^{\rm graph}
    =\int_{D_{AB}}
    \Phi_{\hat d}\!\big(g_{AB}(\xi)\big)
    w_{AB}(\xi)\,\dd \xi,
    \label{eq:graph_energy}
\end{equation}
where $\Phi_{\hat d}$ is a penalty, barrier, or augmented-Lagrangian integrand. The draft implementation supports both contact-sample output and quadrature-based energy output. A downstream finite element solver may use the resulting residuals and tangents, while a collision-only application may use the minimizers and contact normals.

\subsection{Local minimization in the graph domain}

For contact detection, one solves
\begin{equation}
    \xi^*=\argmin_{\xi\in D_{AB}} g_{AB}(\xi).
    \label{eq:xi_min}
\end{equation}
The planned implementation approximates $D_{AB}$ by the overlap of projected reference triangles for the first prototype. A more exact version clips projected Bezier edge curves or uses adaptive sampling. Newton or safeguarded quasi-Newton iterations minimize $g_{AB}$ inside the overlap; active-set logic handles overlap edges. If the minimum is below $\hat d$, a contact sample is emitted:
\begin{equation}
    c=(x_A,x_B,n,g,w,p_A,p_B,\operatorname{type},\operatorname{certificate} ),
\end{equation}
where $x_A=X_A(p_A(\xi^*))$, $x_B=X_B(p_B(\xi^*))$, $g=g_{AB}(\xi^*)$, and $n$ is either $n_c$ or the normalized smooth opposing normal
\begin{equation}
    n(\xi)=\frac{N_A(p_A(\xi))-N_B(p_B(\xi))}{\|N_A(p_A(\xi))-N_B(p_B(\xi))\|}.
\end{equation}
For quadrature, $D_{AB}$ is subdivided adaptively and quadrature points are retained only where $g<\hat d$ or where interval bounds cannot exclude activation.

\subsection{Generic curved patch-patch fallback}
\label{subsec:patch_pair_fallback}

If graphability fails or the graph minimization is uncertain, the algorithm solves the generic curved patch-patch problem
\begin{equation}
    \min_{p\in U_A,\;q\in U_B}\frac{1}{2}\|X_A(p)-X_B(q)\|^2.
    \label{eq:patch_patch}
\end{equation}
Let $r=X_A(p)-X_B(q)$ and $J_A=DX_A(p)$, $J_B=DX_B(q)$. Interior stationary points satisfy
\begin{equation}
    J_A^T r=0,
    \qquad
    J_B^T r=0.
    \label{eq:stationarity_patch}
\end{equation}
A Gauss-Newton step solves
\begin{equation}
    \begin{bmatrix}
    J_A^TJ_A & -J_A^TJ_B\\
    -J_B^TJ_A & J_B^TJ_B
    \end{bmatrix}
    \begin{bmatrix}\Delta p\\ \Delta q\end{bmatrix}
    =-
    \begin{bmatrix}J_A^T r\\ -J_B^T r\end{bmatrix}.
    \label{eq:gn_patch}
\end{equation}
The parameters are projected back to the reference triangles. If a minimizer lies on the boundary of one patch, the local active set is reduced to a curve-surface problem; if it lies on both boundaries, it becomes a curve-curve problem. Vertex cases are handled as further active-set degeneracies. Thus open surface boundaries and sharp features need not be globally pre-classified for the algorithm to remain valid; they appear naturally through the active set. Feature-aware normal splitting can still be used to prevent artificial smoothing across known sharp creases.

\subsection{Conservative fallback and certificates}

A final fallback prevents uncertain results from being silently accepted. For Bezier primitives, the convex hull property gives conservative bounds on spatial extent. Interval ranges or Bernstein bounds on the gap can reject subdomains. Interval Newton or local subdivision is invoked when the standard Newton residual remains large, when the Hessian is ill-conditioned, or when the gap interval straddles the activation threshold. The final emergency fallback is the original linear triangle distance or CCD query with inflated tolerances. This hierarchy is summarized as
\begin{equation}
    \text{BVH} \rightarrow
    \text{linear reject} \rightarrow
    \text{graph certificate} \rightarrow
    \text{2D graph solve} \rightarrow
    \text{4D patch solve} \rightarrow
    \text{certified subdivision}.
\end{equation}

\subsection{Error and convergence estimates}

Let $g$ be the true local signed surface gap and $\hat g$ be the gap computed from curved jets. In a regular contact region with bounded curvature, if both position approximations satisfy $\|\hat X_i-X_i\|\le \eps_i^x$ and normal errors satisfy $\angle(\hat N_i,N_i)\le \eps_i^N$, then the gap error obeys the estimate
\begin{equation}
    |\hat g-g|\le
    \eps_A^x+\eps_B^x
    + C d (\eps_A^N+\eps_B^N)
    + C_\kappa \left((\eps_A^x)^2+(\eps_B^x)^2\right),
    \label{eq:gap_error}
\end{equation}
where $d$ is the local separation scale and $C_\kappa$ depends on curvature bounds. Near contact, $d\approx 0$, so the dominant term is the position approximation error. Under Assumption~\ref{assump:jet_accuracy}, the expected gap and normal errors are
\begin{equation}
    |\hat g-g|=O(h^3),
    \qquad
    \angle(\hat n,n)=O(h^2),
    \label{eq:expected_convergence}
\end{equation}
for smooth regions. In contrast, a globally piecewise-linear surface typically produces lower-order normal behavior and discontinuities across element boundaries. The numerical examples in Sec.~\ref{sec:numerical_design} are designed to measure these rates rather than assume them.

\subsection{Algorithmic summary}

\begin{algorithm}[t]
\caption{Preprocessing for \method{}}
\label{alg:preprocess}
\begin{algorithmic}[1]
\Require Triangulated surface mesh $S$, activation distance $\hat d$, optional shell radius $r$
\Ensure Contact object with curved primitives and inflated BVH
\For{each triangle $T_i\subset S$}
    \State Fit curved jet primitive $\Pc_i=(X_i,N_i,\kappa_i^{\max},\eps_i^x,\eps_i^N,\Cc_i^N)$
    \State Compute or bound $\operatorname{AABB}(X_i(U_i))$
    \State Inflate using Eq.~\eqref{eq:inflated_aabb}
\EndFor
\State Build a 3D BVH over inflated primitive boxes
\State Return primitives, BVH, and optional feature tags
\end{algorithmic}
\end{algorithm}

\begin{algorithm}[t]
\caption{Pairwise contact query}
\label{alg:pair_query}
\begin{algorithmic}[1]
\Require Candidate primitive pair $(\Pc_A,\Pc_B)$
\Ensure Contact samples or patch quadrature data
\If{linear distance rejects pair with conservative error bounds}
    \State \Return empty set
\EndIf
\State Construct contact frame $(t_1,t_2,n_c,x_c)$ using Sec.~\ref{subsec:contact_frame}
\If{normal-cone graphability test in Sec.~\ref{subsec:certificate} passes}
    \State Solve the two-dimensional graph gap problem in Sec.~\ref{subsec:graph_gap}
    \If{solution is certified}
        \State \Return graph contact samples or graph patch quadrature
    \EndIf
\EndIf
\State Solve generic curved patch-patch problem Eq.~\eqref{eq:patch_patch}
\If{patch solve is uncertain}
    \State Invoke interval/Bernstein/subdivision fallback
\EndIf
\State \Return certified contacts or empty set
\end{algorithmic}
\end{algorithm}

\begin{algorithm}[t]
\caption{Global contact detection}
\label{alg:global}
\begin{algorithmic}[1]
\Require Contact objects $A,B$
\Ensure Contact set $\mathcal{K}$
\State Refit or rebuild BVHs if geometry changed
\State $\mathcal{P}\gets$ BVH candidate pairs between $A$ and $B$
\State $\mathcal{K}\gets\emptyset$
\For{each $(\Pc_A,\Pc_B)\in\mathcal{P}$}
    \State $\mathcal{K}\gets \mathcal{K}\cup \textsc{PairwiseContact}(\Pc_A,\Pc_B)$
\EndFor
\State Merge duplicate samples generated by neighboring primitives
\State Cluster graph samples into contact patches if patch output is requested
\State \Return $\mathcal{K}$
\end{algorithmic}
\end{algorithm}

\subsection{Handling open surfaces, closed surfaces, and sharp features}

The method treats open and closed surfaces uniformly at the search level. Closed surfaces may provide a globally consistent outward normal, but this is not required for building the contact frame. Open surfaces do not have a global signed distance, yet the local gap Eq.~\eqref{eq:graph_gap} remains well-defined whenever graphability holds. Surface boundaries are handled by active-set reduction in the 4D fallback or by clipping of the projected graph domain. If the input contains intentional sharp features, normal splitting is recommended: adjacent primitives on opposite sides of a sharp crease should use separate normal cones and should not be smoothed across the crease. If no feature information is available, large normal cones will naturally reduce graphability and send the pair to the generic fallback.

\subsection{Complexity model}

Let $N$ be the number of input triangles, $K_{\rm cand}$ the number of conservative BVH candidate pairs, $K_{\rm graph}$ the number passing graphability, and $K_{\rm fail}$ the number requiring generic or certified fallback. The expected runtime per step has the form
\begin{equation}
    T = O(N\log N) + K_{\rm cand} C_{\rm cert}
        + K_{\rm graph} C_{\rm 2D}
        + K_{\rm fail} C_{\rm fallback},
    \label{eq:complexity}
\end{equation}
where $C_{\rm cert}$ is the inexpensive normal-cone and linear-rejection cost, $C_{\rm 2D}$ is the graph gap minimization or quadrature cost, and $C_{\rm fallback}$ is the cost of generic patch-patch and certified methods. The method is advantageous when regular surface-to-surface contact dominates, so that
\begin{equation}
    K_{\rm fail}\ll K_{\rm graph}\le K_{\rm cand}.
\end{equation}
This is expected in sliding, conforming, and smooth shell contact, but not necessarily in highly interlocking corner collisions. The planned numerical examples explicitly report graph fast-path ratios and fallback rates.

\section{Numerical examples: design and validation protocol}
\label{sec:numerical_design}

This section specifies the numerical examples completed by the current implementation. The tables below are filled from CSV outputs to avoid manual fabrication. The goals remain to measure accuracy, efficiency, robustness, and the practical frequency of the graphability fast path.

\subsection{Implementation variants to compare}

All benchmarks should compare the following variants.
\begin{enumerate}[leftmargin=*]
    \item \textbf{PL-linear.} Standard piecewise-linear triangle-triangle distance and contact normals.
    \item \textbf{PL-refined-$k$.} Uniformly subdivided piecewise-linear surfaces with $k\in\{1,2,3\}$ refinement levels.
    \item \textbf{CJ-4D.} Curved jet primitives with only the generic 4D patch-patch closest-point solve.
    \item \textbf{CJ-Graph.} Full \method{} using normal-cone graphability, 2D graph solve, 4D fallback, and conservative error bounds.
    \item \textbf{CJ-Graph-Patch.} Full graph formulation with contact patch quadrature and symmetric area weighting.
\end{enumerate}

\subsection{Metrics}

For each benchmark, record the following metrics.
\begin{enumerate}[leftmargin=*]
    \item Gap error: $|\hat g-g_{\rm ref}|$ against analytic or highly refined reference geometry.
    \item Normal error: $\angle(\hat n,n_{\rm ref})$.
    \item Contact point error: Euclidean error in closest or active contact locations.
    \item Force smoothness: total variation of contact normal or contact force during sliding.
    \item Runtime breakdown: broad phase, certificate tests, 2D graph solve, 4D fallback, conservative subdivision, and merging.
    \item Candidate statistics: $K_{\rm cand}$, $K_{\rm graph}$, $K_{\rm fail}$, graph fast-path ratio, and fallback reason distribution.
    \item Memory footprint: primitive storage, BVH storage, and optional patch quadrature data.
    \item Robustness: number of Newton failures, interval fallbacks, and missed contacts relative to a conservative reference.
\end{enumerate}

\subsection{Example 1: analytic sphere-plane and sphere-sphere contact}

\textbf{Purpose.} Verify gap and normal accuracy against analytic solutions. A sphere represented by increasingly coarse triangulations contacts a plane and another sphere. The analytic gap, normal, and contact point are known.

\textbf{Protocol.}
\begin{enumerate}[leftmargin=*]
    \item Generate sphere meshes at resolutions $h_1>h_2>\cdots$.
    \item Fit curved jet primitives from the same vertices and vertex normals.
    \item Move the sphere toward the plane in prescribed separations $g\in\{10^{-1},10^{-2},10^{-3}\}R$.
    \item Repeat for sphere-sphere contact with different radii.
\end{enumerate}

\textbf{Expected figures.} Log-log convergence plots for gap and normal errors; runtime versus resolution; graphability ratio versus resolution.

\textbf{Expected conclusion to verify.} The curved-jet graph method should approach high-order gap and normal convergence without increasing primitive count, while PL-refined variants require increasing primitive counts to reach similar accuracy.

\subsection{Example 2: sliding contact and force-normal continuity}

\textbf{Purpose.} Test whether smooth curved jets reduce normal and force jumps during sliding. A smooth ellipsoid or torus patch slides over a curved rigid guide.

\textbf{Protocol.}
\begin{enumerate}[leftmargin=*]
    \item Prescribe a trajectory that makes the contact point cross many original triangle edges.
    \item Use a fixed activation distance and a smooth penalty potential.
    \item Record contact normal, tangential basis, penalty force, and graph contact frame at each time step.
    \item Compare PL-linear, PL-refined, CJ-4D, and CJ-Graph.
\end{enumerate}

\textbf{Metrics.} Force total variation, maximum normal jump between consecutive steps, runtime, and fallback ratio.

\subsection{Example 3: graphability stress test}

\textbf{Purpose.} Quantify where the graph fast path is valid and where fallback is necessary. Construct controlled contacts with varying contact angle and curvature.

\textbf{Protocol.}
\begin{enumerate}[leftmargin=*]
    \item Contact two sinusoidal surfaces with controllable amplitude and wavelength.
    \item Vary normal-cone threshold $\mu\in\{0.2,0.4,0.6,0.8\}$.
    \item Vary local curvature and tangential offset.
    \item Compare 2D graph results against full 4D patch-patch solves.
\end{enumerate}

\textbf{Metrics.} Graph pass ratio, false reject ratio, gap difference between 2D and 4D solves, and condition number versus $1/\mu$.

\subsection{Example 4: open shell boundary contact}

\textbf{Purpose.} Demonstrate that the method handles open surfaces without a global signed distance or manually specified cut. A thin shell strip slides against a rigid curved surface; contact may occur in the interior and near the open boundary.

\textbf{Protocol.}
\begin{enumerate}[leftmargin=*]
    \item Generate an open rectangular shell with smooth curvature.
    \item Move it over a cylinder or saddle surface.
    \item Track active-set transitions from surface-surface to curve-surface and curve-curve fallback.
\end{enumerate}

\textbf{Metrics.} Number of boundary contacts, missed-contact checks against refined PL reference, and runtime overhead of active-set fallback.

\subsection{Example 5: contact patch quadrature and pressure distribution}

\textbf{Purpose.} Validate the patch-integral formulation. Compress a smooth elastic-like surface patch against a rigid indenter using prescribed displacements, without solving a full bulk FE problem in the first version.

\textbf{Protocol.}
\begin{enumerate}[leftmargin=*]
    \item Use a paraboloid indenter and a smooth deformable surface represented kinematically.
    \item Evaluate $E_{AB}^{\rm graph}$ using adaptive quadrature on $D_{AB}$.
    \item Compare total normal force and contact area against analytic Hertzian trends where applicable or against a very fine quadrature reference.
\end{enumerate}

\textbf{Metrics.} Contact area error, total force error, quadrature count, and symmetry difference between $A\rightarrow B$ and $B\rightarrow A$.

\subsection{Example 6: integration with a simple nonlinear solver}

\textbf{Purpose.} Show that the contact front end can drive a small mechanical solve. This is not intended to replace a production FE solver, but to verify API design and differentiability.

\textbf{Protocol.}
\begin{enumerate}[leftmargin=*]
    \item Use a small mass-spring or shell-like surface model with prescribed displacement.
    \item Use the graph patch energy or contact samples with a penalty/barrier potential.
    \item Run quasi-static Newton iterations and compare convergence against PL-linear contact.
\end{enumerate}

\textbf{Metrics.} Newton iterations, residual norms, contact energy smoothness, and robustness under sliding.

\subsection{Example 7: scalability on random smooth surfaces}

\textbf{Purpose.} Measure broad-phase and narrow-phase scaling independent of analytic geometry.

\textbf{Protocol.}
\begin{enumerate}[leftmargin=*]
    \item Generate random smooth height-field surfaces, remeshed at multiple resolutions.
    \item Bring pairs into near contact with random offsets and rotations.
    \item Record all runtime and candidate statistics.
\end{enumerate}

\textbf{Metrics.} Runtime versus $N$, active pair count, graph pass ratio, memory usage, and parallel efficiency if the backend is threaded.

\subsection{Example 8: reserved industrial benchmark}

A complex industrialized scenario is intentionally reserved until Examples 1--7 are complete and stable. The implementation plan leaves this scenario as a separate stage. Candidate scenarios include gear-tooth flank contact, bearing roller-race contact, seal lip contact on a shaft, or turbine blade tip rub. The final selection should be based on available geometry, reference data, and expected CMAME relevance. No results should be reported for this scenario until the validation sequence above passes.

\subsection{Planned tables}

\subsection{Prototype snapshot from direct CSV outputs}

\paragraph{Example 1 (analytic and convergence checks).}
The direct CSV rows are shown in Table~\ref{tab:accuracy_snapshot}.

\begin{table}[h]
\centering
\caption{Example 1 prototype snapshot.}
\label{tab:accuracy_snapshot}
\begin{tabular}{l l r r r r}
\toprule
Case & Method & Mesh size & Gap proxy & Normal error proxy & Runtime (ms) \\
\midrule
Sphere--plane (sep 0.01) & CJ-Graph & 6 & $1.0\times 10^{-2}$ & N/A & 40.8651 \\
Sphere--sphere (sep 0.01) & CJ-Graph & 6 & $7.8666\times 10^{-3}$ & N/A & 2401.5894 \\
Sphere--plane (sep 0.03) & CJ-Graph & 6 & $3.0\times 10^{-2}$ & N/A & 40.4306 \\
Sphere--sphere (sep 0.03) & CJ-Graph & 6 & $7.8666\times 10^{-3}$ & N/A & 2434.9042 \\
Sphere--plane (sep 0.01) & CJ-Graph & 10 & $1.0\times 10^{-2}$ & N/A & 123.2006 \\
Sphere--sphere (sep 0.01) & CJ-Graph & 10 & $6.4667\times 10^{-4}$ & N/A & 21663.5375 \\
Sphere--plane (sep 0.03) & CJ-Graph & 10 & $3.0\times 10^{-2}$ & N/A & 121.9446 \\
Sphere--sphere (sep 0.03) & CJ-Graph & 10 & $6.4667\times 10^{-4}$ & N/A & 21544.1193 \\
\bottomrule
\end{tabular}
\end{table}

\paragraph{Examples 2--7 (prototype status snapshot).}

\begin{table}[h]
\centering
\caption{Prototype snapshot for graphability, fallback, and scaling diagnostics.}
\label{tab:graph_snapshot}
\begin{tabular}{l r r r r r}
\toprule
Example & Candidate pairs & Graph pass & 4D fallback & Certified fallback & Fast-path ratio \\
\midrule
Example 2 (sliding) & 2448 & 0 & 0 & 2448 & 0.000 \\
Example 3 (graphability stress) & 0 & 0 & 0 & 0 & 0.000 \\
Example 4 (open shell) & 0 & 0 & 0 & 0 & 0.000 \\
Example 5 (patch/symmetry) & 0 & 0 & 0 & 0 & 0.000 \\
Example 6 (solver coupling) & 0 & 0 & 0 & 0 & 0.000 \\
Example 7 (res 10) & 40000 & 0 & 2557 & 40000 & 0.000 \\
Example 7 (res 14) & 153664 & 0 & 5902 & 153664 & 0.000 \\
\bottomrule
\end{tabular}
\end{table}

\paragraph{Remark on missing physics outputs.}
Example 2, 3, 4, 5, and 6 currently report $inf$/zero gaps and zero graph candidates in several rows, consistent with the present conservative settings and reduced query geometry. Example 7 is currently the only scalable example that produces finite best gaps under this run.

\section*{Acknowledgements}
This draft intentionally leaves funding, author contributions, and numerical results blank.

\begin{thebibliography}{99}

\bibitem{laursen2002}
T. A. Laursen, \emph{Computational Contact and Impact Mechanics: Fundamentals of Modeling Interfacial Phenomena in Nonlinear Finite Element Analysis}, Springer, Berlin, 2002.

\bibitem{wriggers2006}
P. Wriggers, \emph{Computational Contact Mechanics}, 2nd ed., Springer, Berlin, 2006.

\bibitem{puso2004mortar}
M. A. Puso and T. A. Laursen, A mortar segment-to-segment contact method for large deformation solid mechanics, \emph{Computer Methods in Applied Mechanics and Engineering} 193 (2004) 601--629.

\bibitem{puso2004friction}
M. A. Puso and T. A. Laursen, A mortar segment-to-segment frictional contact method for large deformations, \emph{Computer Methods in Applied Mechanics and Engineering} 193 (2004) 4891--4913.

\bibitem{sahu2025midplane}
I. Sahu and N. Petrinic, Midplane based 3D single pass unbiased segment-to-segment contact interaction using penalty method, \emph{Computer Methods in Applied Mechanics and Engineering} 447 (2025) 118335.

\bibitem{li2020ipc}
M. Li, Z. Ferguson, T. Schneider, T. Langlois, D. Zorin, D. Panozzo, C. Jiang, and D. M. Kaufman, Incremental Potential Contact: Intersection- and inversion-free, large-deformation dynamics, \emph{ACM Transactions on Graphics} 39(4) (2020) Article 49.

\bibitem{huang2025gcp}
Z. Huang, M. Paik, Z. Ferguson, D. Panozzo, and D. Zorin, Geometric Contact Potential, \emph{ACM Transactions on Graphics} 44(4) (2025).

\bibitem{du2024smooth}
Y. Du, Y. Li, S. Coros, and B. Thomaszewski, Robust and artefact-free deformable contact with smooth surface representations, \emph{Computer Graphics Forum} (2024).

\bibitem{ferguson2023highorder}
Z. Ferguson, P. Jain, D. Zorin, T. Schneider, and D. Panozzo, High-order Incremental Potential Contact for elastodynamic simulation on curved meshes, \emph{ACM Transactions on Graphics} 42(4) (2023).

\bibitem{chen2024tdib}
X. Chen, C. Yu, X. Ni, M. Chu, B. Wang, and B. Chen, A time-dependent inclusion-based method for continuous collision detection between parametric surfaces, \emph{ACM Transactions on Graphics} 43(6) (2024).

\bibitem{chen2025ogc}
A. H. Chen, J. Hsu, Z. Liu, M. Macklin, Y. Yang, and C. Yuksel, Offset Geometric Contact, \emph{ACM Transactions on Graphics} 44(4) (2025).

\bibitem{pelletier2025trisdf}
J. Pelletier-Guenette and S. Andrews, Real-time triangle-SDF continuous collision detection, \emph{Proceedings of the ACM SIGGRAPH / Eurographics Symposium on Computer Animation} (2025).

\bibitem{sawhney2018bff}
R. Sawhney and K. Crane, Boundary First Flattening, \emph{ACM Transactions on Graphics} 37(1) (2018) Article 5.

\bibitem{vlachos2001pn}
A. Vlachos, J. Peters, C. Boyd, and J. L. Mitchell, Curved PN triangles, \emph{Proceedings of the Symposium on Interactive 3D Graphics} (2001) 159--166.

\end{thebibliography}

\end{document}

